\section{Introduction}
In the past year, language models such as Code Llama \citep{roziere2023code}, DeepSeek-Coder \citep{guo2024deepseek}, and GPT-4 \citep{openai2023gpt} have demonstrated great advances in programming and software engineering tasks. Their success has primarily been due to their strong code generation abilities, as measured by benchmarks such as HumanEval \citep{chen2021evaluating} and MBPP \citep{austin2021program}, as well as their usefulness in general purpose code-writing. However, being proficient at coding is a complex skill that involves more than just code generation: when a skilled programmer writes a piece of code, they generally have a multifaceted understanding of their code: they can generally go back and check if it is correct, reason about the execution behaviour of the code, and debug the code if there are errors. On the other hand, while dozens of benchmarks have shown that these models can generate code, it is unclear whether language models have this same level of understanding of that code.

In textual domains, \citet{west2023generative} discovered ``the generative AI paradox'': language models may have strong generative capabilities while failing to understand what they generate. We investigate this finding from the perspective of coding: \textit{to what extent do code language models understand the programs that they write}? 

The advantage of exploring this question through code is that programs have a precise semantics, so code understanding can be analyzed by investigating whether models can reason through tasks requiring a nuanced understanding of these semantics. We probe code understanding through the lens of three tasks: 1) correctness checking, 2) execution prediction, and 3) program repair. In the first task, the model is asked to assess whether a short piece of code correctly implements a natural language specification (sometimes with test cases). In the second task, the model is given a program-input pair and asked to predict the output of executing the program on the given input. For this task to be fair, we ensure the programs are generally short and that execution does not require complex calculation. In the final task, the model is given an incorrect program alongside its original specification and asked to correct it. 

While a similar task of re-ranking methods have been employed to improve code generation \citep{inala2022fault, ni2023lever}, this work is (to our knowledge) the first to study the correctness checking task in a zero-shot setting for general-purpose coding models.\footnote{A concurrent work \citep{singhal2024nofuneval} considers a related task of identifying which one of two programs has a bug, and incorrect programs were generated by having GPT-4 synthetically introduce bugs.} When writing programs following a well-formulated specification, correctness checking is a foundational prerequisite to code understanding: how can someone understand a program if they can't even determine if it performs the intended functionality? In addition, self-repair relies on the model knowing that the code is wrong in the first place. The second and third tasks were introduced in \citep{austin2021program, gu2024cruxeval} and \citep{chen2024teaching, olausson2024repair}, and we employ them to draw insights for code understanding.

Our contributions are as follows:
\begin{itemize}
    \item First, we conduct an evaluation and analysis of the correctness-checking capabilities of six code LMs from the StarCoder, DeepSeek, and CodeLlama families on three datasets: HumanEval, LeetCode, and ODEX. We find that open models score between random guessing and 70\%. Through our analysis, we also discover that models are generally biased into judging programs as ``Correct'' but that stronger models have decently calibrated log-probabilities even without CoT.
    \item Second, we surprisingly observe no correlation between problem difficulty and correctness checking ability. However, models are better at executing and repairing programs for easier problems. Nevertheless, we observe that models frequently cannot necessarily check, execute, and repair solutions for problems that it can solve over 80\% of the time, suggesting that \textit{models do not have a good understanding of programs that they write}. 
    \item Third, we provide a detailed qualitative analysis of GPT-4 on correctness checking, shedding light on the failure modes and potential areas of improvement for future models.
\end{itemize}